\chapter{Combinatorics}

\section{Hoán vị}
	\subsection{Giai thừa}
		\import{factorial.tex}

	\subsection{Hoán vị không điểm cố định (số lệch)}
		Hoán vị của một tập sao cho không phần tử nào nằm ở vị trí ban đầu. Ký hiệu $\lfloor x \rceil$ là làm tròn đến số nguyên gần nhất.
		\[ \mkern-2mu D(n) = (n-1)(D(n-1)+D(n-2)) = n D(n-1)+(-1)^n = \left\lfloor\frac{n!}{e}\right\rceil \]

	\subsection{Bổ đề Burnside}
		Cho nhóm các phép đối xứng $G$ tác động lên tập $X$. Số phần tử của $X$ \emph{xét tới đối xứng} bằng
		 \[ {\frac {1}{|G|}}\sum _{{g\in G}}|X^{g}|, \]
		 trong đó $X^{g}$ là tập các phần tử bị $g$ giữ cố định ($g.x = x$).

		 Nếu $f(n)$ đếm số cấu hình có độ dài $n$, bỏ qua đối xứng quay ($G = \mathbb Z_n$):
		 \[ g(n) = \frac 1 n \sum_{k=0}^{n-1}{f(\text{gcd}(n, k))} = \frac 1 n \sum_{k|n}{f(k)\phi(n/k)}. \]

\section{Phân hoạch}
	\subsection{Hàm phân hoạch số}
		Số cách viết $n$ thành tổng các số nguyên dương, không tính thứ tự.
		\[ p(0) = 1,\ p(n) = \sum_{k \in \mathbb Z \setminus \{0\}}{(-1)^{k+1} p(n - k(3k-1) / 2)} \]

		\begin{center}
		\begin{tabular}{c|c@{\ }c@{\ }c@{\ }c@{\ }c@{\ }c@{\ }c@{\ }c@{\ }c@{\ }c@{\ }c@{\ }c@{\ }c}
			$n$    & 0 & 1 & 2 & 3 & 4 & 5 & 6  & 7  & 8  & 9  & 20  & 50  & 100 \\ \hline
			$p(n)$ & 1 & 1 & 2 & 3 & 5 & 7 & 11 & 15 & 22 & 30 & 627 & $\mathtt{\sim}$2e5 & $\mathtt{\sim}$2e8 \\
		\end{tabular}
		\end{center}

\section{Các dãy số thường gặp}
	\subsection{Số Stirling loại 1}
		Số hoán vị của $n$ phần tử có đúng $k$ chu trình.
		\[ c(n,k) = c(n-1,k-1) + (n-1) c(n-1,k),\ c(0,0) = 1 \]
		\[ \textstyle \sum_{k=0}^n c(n,k)x^k = x(x+1) \dots (x+n-1) \]

	\subsection{Số Stirling loại 2}
		Số cách phân hoạch $n$ phần tử phân biệt thành đúng $k$ nhóm không phân biệt thứ tự.
		\[ S(n,k) = S(n-1,k-1) + k S(n-1,k),\ S(n,1) = S(n,n) = 1 \]
		\[ S(n,k) = \frac{1}{k!}\sum_{j=0}^k (-1)^{k-j}\binom{k}{j}j^n \]

	\subsection{Số Bell}
		Tổng số cách phân hoạch $n$ phần tử phân biệt: $B(n) = \sum_k S(n,k)$.
		$B(n) = 1, 1, 2, 5, 15, 52, 203, 877, 4140, 21147, \dots$

	\subsection{Cây không gốc có gán nhãn}
		Số cây trên $n$ đỉnh: $n^{n-2}$ \\
		Số cách nối $k$ cây có sẵn (kích thước $n_i$) thành một cây: $n_1n_2\cdots n_k n^{k-2}$ \\
		Số cây với bậc các đỉnh là $d_i$: $(n-2)! / ((d_1-1)! \cdots (d_n-1)!)$

	\subsection{Số Catalan}
		\[ C_n=\frac{1}{n+1}\binom{2n}{n}= \binom{2n}{n}-\binom{2n}{n+1} = \frac{(2n)!}{(n+1)!n!} \]
		\[ C_0=1,\ C_{n+1} = \frac{2(2n+1)}{n+2}C_n,\ C_{n+1}=\sum C_iC_{n-i} \]
		${C_n = 1, 1, 2, 5, 14, 42, 132, 429, 1430, 4862, 16796, 58786, \dots}$
		\begin{itemize}
		\item[-] Số đường đi đơn điệu nằm dưới đường chéo trong lưới $n\times n$.
		\item[-] Số dãy ngoặc đúng gồm $n$ cặp ngoặc.
		\item[-] Số cây nhị phân có $n+1$ lá, mỗi nút có 0 hoặc 2 con.
		\item[-] Số cây có thứ tự gồm $n+1$ đỉnh.
		\item[-] Số cách tam giác hóa đa giác lồi $n+2$ cạnh.
		\item[-] Số hoán vị của $[n]$ không chứa dãy con tăng độ dài 3.
		\end{itemize}
